\section{Discussion}
\label{sec:discussion}

\subsection{Why behavioural detectors fail against adaptive attackers}
The graph networks learn exactly what they are trained to learn: in the observed campaigns, malicious delegations are submitted by sponsors in batches, from accounts funded by earlier drains, to contracts whose earlier victims have already been swept. These regularities are statistically strong---they yield an AUROC of $0.999$ on unseen families---but they are properties of the attackers' current \emph{operations}, not of theft. Each can be removed for the cost of a few transactions: a leaked-key attacker can sign and submit from the victim's account, fund gas through a fresh path, or deploy one contract per victim. Our suite shows that such changes leave the ranking of graph models largely intact, yet move malicious scores below any threshold calibrated on historical behaviour. The general lesson for security applications of temporal graph learning is that high accuracy on behaviour-rich historical data says little about robustness when the behaviour is attacker-chosen. Evaluations should include attacks on the features a model actually relies on.

\subsection{Why canonical code patterns are robust}
A theft mechanism must exist in the delegate's code: value has to be forwarded to an attacker-controlled destination, or an attacker-controlled caller has to be granted arbitrary execution. The cheapest way to field a new campaign is to copy an existing mechanism and vary only its surface---addresses, filler instructions, equivalent opcodes, appended data. \dn{}'s canonicalisation is designed around precisely this economics. It removes, by construction, every transformation in the class $\mathcal{M}$, so copying no longer evades detection, and an attacker who wants to escape the memory must implement theft differently, which is substantially more expensive than re-deploying a template. The memory's family exclusion ensures that what \dn{} learns is resemblance to \emph{other} known mechanisms, which is the property needed for new families. That first-victim detection is \dn{}'s strongest accuracy slice is consistent with this interpretation.

\subsection{Deployment}
\paragraph{Wallets} \dn{} needs only the delegate address of a signature request: the wallet retrieves the code, canonicalises it and scores it against a memory distributed as a list of four-gram sets. At a 1\% validation false-positive operating point it warns on $98.4\%$ of malicious requests while interrupting about one in $600$ benign ones, and each warning cites the specific known contract the request resembles.
\paragraph{Monitoring services} Explorers and security providers can run \dn{} alongside behavioural detectors. The latter are more accurate once a campaign has accumulated victims, while \dn{} covers the first victims, disguised re-deployments and attackers who abandon sponsored submission. A disagreement between the two is itself informative: a high code score with a low behavioural score is the signature of an attacker who has changed operations but not mechanism.
\paragraph{Memory maintenance} New confirmed families are added to the memory without retraining. Because labels can be derived automatically from outcomes (Section~\ref{sec:data}), the memory can be refreshed continuously, with human review reserved for borderline contracts.

\subsection{What an attacker can still do}
\dn{}'s guarantee covers the transformation class $\mathcal{M}$. Attackers could move beyond it, at higher cost:
\begin{itemize}
\item \emph{Semantic obfuscation}: opaque predicates, computed jump targets, or arithmetic encoding of the attacker address. Computed jump targets defeat static reachability; this is a known limit of static analysis that symbolic execution of the dispatcher would address.
\item \emph{Indirection}: a minimal proxy that \texttt{DELEGATECALL}s into separately deployed logic. Our idioms flag delegate calls to calldata- or storage-derived targets, and resolving the implementation and canonicalising it as well is a direct extension.
\item \emph{Novel mechanisms}: a genuinely new theft mechanism has no close memory match and receives an intermediate score. This is the regime where behavioural evidence, accumulated over the first victims, becomes valuable.
\end{itemize}
Each of these raises the attacker's cost relative to today's practice of re-deploying copied templates.

\subsection{Limitations}
\paragraph{Benign coverage} Our labelling rules require at least five delegators and sustained owner activity before a contract is labelled benign, which yields $11$ benign contracts in $9$ families. This is conservative---it avoids labelling as benign a contract whose victims simply have not been drained yet---but it bounds the precision of contract-level estimates, as the bootstrap intervals of Table~\ref{tab:sig} show. Curated lists of official wallet implementations would enlarge the benign set.
\paragraph{Assets and chains} Activity was collected for native ether, including internal transfers, but not for ERC-20 tokens, because free endpoints do not serve historical logs at the required scale. Drainers that steal only tokens are therefore labelled from their ether behaviour alone. The study covers Ethereum; other chains that adopted EIP-7702 lack free historical trace access, and per-address explorer data would enable a cross-chain evaluation.
\paragraph{Sampling} Activity was collected for a per-contract stratified sample of up to $1{,}000$ delegators. The full authorization scan is released, so that complete histories can be added with commercial data access.
\paragraph{Window} The window covers the first two months after activation, when sweepers dominated. As the mechanism matures, the balance between benign and malicious delegations, and the attackers' operations, will change. Continuous re-evaluation under the zero-day protocol is the appropriate response.

\subsection{Threats to validity}
\emph{Internal}: labels are derived from outcomes with fixed rules, and mislabelled contracts would affect every method equally. The rules are reported in full, the evidence behind every label is released, and the independent cross-check found no benign contract associated with a known scam address. \emph{Construct}: AUROC and fixed-false-positive recall measure ranking and operating-point behaviour respectively, and we report both, together with contract-level aggregation to prevent large campaigns from dominating. \emph{External}: results depend on the attack mix of the observation window; the adaptive-adversary suite probes deviations from it, but only along the dimensions we modelled. \emph{Conclusion}: with $60$ first-victim contracts per seed, first-victim differences between learned methods are within bootstrap uncertainty, and we draw no conclusion from them beyond \dn{}'s consistently high and stable performance on that slice.

\subsection{Ethics and data release}
All data are public ledger records. Addresses are pseudonymous, and we release no off-chain identities, link no addresses to individuals, and publish no information not already present on chain. The dataset documents attacker infrastructure that is already publicly visible, and its release supports defenders more than attackers, who already possess it. The canonicaliser and mutator describe transformations attackers already apply; publishing them strengthens defences without enabling new attacks. The dataset is accompanied by a data card describing collection, labelling rules, known limitations and intended use.
